\section{Marco Teórico}

\subsection{Análisis Paralingüístico}

El análisis paralingüístico estudia los elementos de la comunicación vocal más allá
del contenido lingüístico. Específicamente, analiza el tono (altitud de la voz), el
ritmo (velocidad de habla), el volumen (intensidad sonora) y la calidad vocal
(características específicas de la voz, como si sonara ronca o clara). Estas
características transmiten información sobre el estado emocional del hablante,
independientemente del texto pronunciado. La Tabla~\ref{tab:paralinguistic} resume las
principales características y sus implicaciones emocionales \cite{jha2022, lian2023}.

\begin{table}[!t]
\renewcommand{\arraystretch}{1.3}
\caption{Características Paralingüísticas y sus Implicaciones Emocionales}
\label{tab:paralinguistic}
\centering
\begin{tabular}{p{2cm} p{5.5cm}}
\toprule
\textbf{Característica} & \textbf{Implicaciones Emocionales} \\
\midrule
Pitch (F0)        & Se eleva con emociones de alta activación; desciende con baja activación. \\
Intensidad        & Aumenta con activación emocional y estrés comunicativo. \\
Tasa de habla     & Se acelera con ansiedad; se ralentiza con tristeza o depresión. \\
Jitter / Shimmer  & Variaciones microtemporales que aumentan con tensión emocional y física. \\
MFCCs             & Capturan timbre vocal para caracterización holística del estado emocional. \\
\bottomrule
\end{tabular}
\end{table}

\subsection{Concepto de Línea Base Emocional}

La alta variabilidad interindividual hace necesario establecer una referencia
personalizada para cada hablante. El \textit{baseline} emocional individual es el
estado emocional neutral de esa persona en particular. Los picos se detectan como
desviaciones significativas respecto a esta \textit{baseline}, no respecto a umbrales
absolutos, lo que permite adaptarse a las características únicas de cada voz
\cite{gat2022}.

\subsection{Machine Learning para Detección de Emociones}

\begin{itemize}
    \item \textbf{Detección estadística:} Z-scores y umbrales de desviación respecto
          de la línea base emocional.
    \item \textbf{Métodos clásicos:} Support Vector Machines (SVM), Random Forest y
          Isolation Forest, algoritmos de aprendizaje automático utilizados para la
          detección de anomalías multivariadas \cite{wagner2023, xu2023, nath2024}.
\end{itemize}

\subsection{Revisión de Antecedentes}

\subsubsection{Contexto del Problema en los Call Centers}

Los call centers constituyen entornos de alta carga emocional en los que miles de
interacciones diarias pueden desembocar en picos críticos de ira, frustración o
ansiedad. La evaluación manual sigue siendo el estándar dominante, pero conlleva que
un supervisor dedique tiempo completo a cada grabación. Feng y Devillers (2023)
evidencian que los modelos entrenados con datos espontáneos de call centers superan
consistentemente a los entrenados con habla actuada \cite{feng2023}. El análisis
basado exclusivamente en el texto transcrito resulta insuficiente para capturar
estados emocionales encubiertos, como el sarcasmo o la frustración sutil, que se
manifiestan con claridad en características acústicas como la elevación del pitch, la
aceleración de la tasa de habla o el incremento de las variaciones de jitter y
shimmer \cite{jha2022, lian2023}.

\subsubsection{Características Acústicas para la Detección de Emociones}

La frecuencia fundamental (F0 o pitch), que mide la frecuencia más baja de la voz,
es uno de los correlatos acústicos más estudiados: incrementos sostenidos de F0 se
asocian con emociones de alta activación, como la ira o el miedo, mientras que
descensos de F0 se correlacionan con la tristeza o la calma \cite{jha2022, lian2023}.
Las métricas de calidad vocal jitter y shimmer miden las perturbaciones
microtemporales en la periodicidad del ciclo glótico; jitter se refiere a la
variabilidad en la frecuencia de la voz, y shimmer a la variabilidad en la amplitud.
Zhang et al. (2023) demostraron en IEEE que estas características capturan diferencias
prosódicas significativas bajo una alta carga emocional \cite{zhang2023}. Lian et al.
(2023) confirman que jitter y shimmer, junto con la razón armónico-ruido (HNR),
capturan la aspereza y la ronquera asociadas con el habla ansiosa o de alta activación
emocional \cite{lian2023}. Los MFCC capturan el timbre vocal convirtiendo las señales
acústicas en una representación numérica de cómo el oído humano percibe el espectro
de frecuencia, y han sido ampliamente utilizados como entrada para clasificadores de
emociones \cite{jha2022, abdelhamid2022}. Wagner et al. (2023) demostraron en IEEE
TPAMI que arquitecturas de transformadores combinadas con eGeMAPS logran cerrar la
brecha de valencia en el reconocimiento de emociones en habla espontánea
\cite{wagner2023}.

\subsubsection{Normalización por Hablante y Línea Base Emocional}

La variabilidad interhablante vuelve ineficaces los umbrales absolutos y justifica el
uso de \textit{baselines} personalizadas. Gat et al. (2022) proponen en ICASSP un
método de normalización por hablante para sistemas de SER basados en
autosupervisión, demostrando que la normalización a nivel de hablante mejora
significativamente el desempeño en habla espontánea \cite{gat2022}. Los tres métodos
de \textit{baseline} más estudiados son: (1) \textit{baseline} de inicio —primeros
2--5 minutos de la llamada—; (2) \textit{baseline} por mediana global —más robusta
ante valores atípicos—; y (3) \textit{baseline} adaptativa con ventana deslizante
—actualización continua durante la llamada—. La elección entre estos métodos
constituye uno de los objetivos específicos de esta investigación.

\subsubsection{Algoritmos de Machine Learning para la Detección de Picos}

El enfoque estadístico basado en Z-scores calcula la desviación de cada ventana
temporal respecto a la línea de base del hablante y marca como pico toda observación
que supere un umbral de desviación estándar. Xu et al. (2023) proponen Deep Isolation
Forest en IEEE TKDE, demostrando que extensiones neurales del iForest clásico superan
a los métodos lineales para la detección de anomalías en espacios de alta dimensión,
lo cual es directamente aplicable al espacio multivariado de features acústicas
\cite{xu2023}. Nath et al. (2024) reportan que Random Forest ofrece un equilibrio
favorable entre la interpretabilidad y la precisión en la clasificación de la valencia
emocional \cite{nath2024}. Jha et al. (2022) concluyen que la combinación de
características prosódicas y espectrales es más robusta que cualquiera de ellas por
separado \cite{jha2022}. Gao et al. (2024) presentan en ICASSP un enfoque de
fine-tuning en dos etapas con adaptadores para SER, relevante cuando el corpus de
llamadas reales disponible es limitado \cite{gao2024}.

\subsubsection{Datasets y Validación Experimental}

RAVDESS es un corpus actuado con actores profesionales ampliamente utilizado como
benchmark. Abdelhamid et al. (2022) reportan en IEEE Access que un CNN+LSTM entrenado
sobre RAVDESS alcanza precisiones superiores al 85\,\% en condiciones controladas,
aunque el rendimiento cae al aplicarse a la habla espontánea sin adaptación de dominio
\cite{abdelhamid2022}. IEMOCAP ofrece conversaciones elicitadas con mayor proximidad
al habla natural; Wagner et al. (2023) lo usan como benchmark principal para demostrar
el cierre de la brecha de valencia \cite{wagner2023}. De Lope y Graña (2023) presentan
en Neurocomputing una revisión continua del estado del arte en SER, señalando que la
generalización a dominios espontáneos es el reto abierto más relevante \cite{delope2023}.
Grósz et al. (2022) demuestran en ACM MM que los embeddings de Wav2Vec~2.0 transfieren
eficazmente a tareas paralingüísticas con datos limitados \cite{grosz2022}. Khare et
al. (2024) concluyen en Information Fusion que la combinación de características
acústicas clásicas con representaciones aprendidas es actualmente la estrategia de
mayor rendimiento en escenarios reales \cite{khare2024}.
